\section{Robotics：自动化的是任务，不是责任}

新硬件和模型进入 physical AI 后，会面对医疗、robotaxi 与 harsh-environment robotics。课堂把它们放在同一问题下，但工程成熟度差异巨大。本节的桥接问题是：系统到底替代了哪个 task、保留了哪个 human role、出现 sensor 或 control failure 时怎样安全退出。

\subsection{三类 robotics 的不同 safety case}

Healthcare robot 可能提高 precision，但 surgeon 和 hospital 承担 clinical responsibility；robotaxi 在开放道路持续感知与控制，需要 operational design domain 和远程支援；lunar/water-scarcity rover 运行在通信延迟和极端环境下，强调 autonomy、fault tolerance 与 energy management。三者不能共享一份泛化“robotics benchmark”。

\begin{importantbox}{Physical AI 的最小 safety case}
\begin{enumerate}
  \item 定义 operational design domain 与禁止条件；
  \item 给出 sensor coverage、uncertainty 和 degradation mode；
  \item 说明 human override、safe stop 与 recovery；
  \item 记录 action、state、software/model version 和 incident；
  \item 用真实场景和 rare event 验证，而不仅是平均成功率。
\end{enumerate}
\end{importantbox}

\begin{warningbox}{劳动转移不是自动的社会收益}
把 repetitive task 交给 robot 可能释放人力，也可能造成技能错配、地区不均与监督劳动增加。系统方案需要同时设计 retraining、transition support、liability 和 service access，不能把“更高价值工作”当作自动发生的结果。
\end{warningbox}

\subsection{本章小结}

Robotics 的成熟度由 task boundary、safety case、human role 和 incident recovery 决定。自动化某个动作不等于转移责任，更不等于整个职业或行业已经被自动化。

